\section{拓扑流形}

本节介绍\textbf{拓扑流形}，这是最基本的流形类型. 我们假设读者已经熟悉拓扑空间的定义和基本性质，详见附录A.

设$M$是一个拓扑空间. 若$M$满足以下性质，则称$M$为\textbf{维数}$\pmb{n}$的\textbf{拓扑流形}，或\textbf{拓扑$\pmb{n}$-流形}：
\begin{enumerate}
    \item[(i)] $M$是Hausdorff空间：对任意两个不同的点$p,q\in M$，存在不相交的开子集$U,V\subseteq M$，使得$p\in U$且$q\in V$.
    \item[(ii)] $M$是第二可数的：$M$的拓扑具有\textbf{可数基}.
    \item[(iii)] $M$是局部欧氏的，维数为$\pmb{n}$：$M$的每一点都有一个邻域，该邻域与$\mathbb{R}^n$的某个开子集同胚.
\end{enumerate}

第三条性质更具体地意味着，对每一点$p\in M$，我们都能找到：
\begin{itemize}
    \item $M$的一个包含$p$的开子集$U$；
    \item $\mathbb{R}^n$的一个开子集$\hat{U}$；
    \item 一个同胚$\varphi\colon U\to\hat{U}$.
\end{itemize}

\thmbase{习题 1.1}{
    证明：若在定义中将“$U$与$\mathbb{R}^n$的某个开子集同胚”替换为“$U$与$\mathbb{R}^n$中的某个开球同胚”，或者替换为“$U$与$\mathbb{R}^n$本身同胚”，则得到的是拓扑流形的等价定义.
}

若$M$是拓扑流形，我们常将$M$的维数简记为$\dim M$. 在非正式场合，有时用“设$M^n$是一个流形”作为“设$M$是一个维数为$n$的流形”的简写. 上标$n$并不是流形名称的一部分，通常在首次出现之后便不再包含在记号中.

需要特别注意的是，根据定义，每个拓扑流形都有一个特定且良定义的维数. 因此，我们根本不把混合维数的空间（例如一个平面与一条直线的不交并）视为流形. 在第17章中，我们将利用de Rham上同调理论来证明下述定理，它表明（非空）拓扑流形的维数实际上是一个拓扑不变量.

\thmbase{定理 1.2（维数的拓扑不变性）}{
    非空$n$维拓扑流形不可能与$m$维流形同胚，除非$m=n$.
}

其证明见定理 17.26. 在第2章中，我们还将证明一个相关但较弱的定理（维数的微分同胚不变性，定理 2.17）. 另见 [LeeTM, Chap.~13]，其中利用奇异同调论给出了定理 1.2 的另一种证明.

空集对任意$n$都满足拓扑$n$-流形的定义. 在大多数情况下，我们将忽略这一特例（有时甚至可能忘了说明）. 但由于在某些情形下允许空流形是有用的，我们选择不把它从定义中排除.

拓扑$n$-流形最基本的例子是$\mathbb{R}^n$本身. 它是Hausdorff空间，因为它是度量空间；它是第二可数的，因为所有以有理点为中心、以有理数为半径的开球构成了其拓扑的一个可数基.

要求流形具备这些性质有助于保证其行为符合我们对欧氏空间的经验预期. 例如，容易验证在Hausdorff空间中，有限子集都是闭集，而且收敛序列的极限是唯一的（见附录A中的习题 A.11）. 第二可数性的动机不那么明显，但它在全书中有重要的
\begin{figure}[htbp]
    \centering
    \includegraphics[width=0.6\textwidth, height=7cm, keepaspectratio]{images/deaff15e2ad797c3025490c53a70167ac1f49d2ca6bd26b4e4464def74e207c4.jpg}
    \caption{坐标卡}
\end{figure}
推论，主要源于单位分解的存在性（见第2章）.

在实践中，Hausdorff性质和第二可数性通常都很容易验证，特别是对于那些由其他流形构造出来的空间，因为这两条性质都被子空间和有限乘积所继承（命题 A.17 和 A.23）. 特别地，拓扑$n$-流形的每个开子集本身也是一个拓扑$n$-流形（当然要赋予子空间拓扑）.

应当指出，有些作者选择从流形的定义中去掉Hausdorff性质或第二可数性，或者两者都去掉. 然而，关于流形的大多数有趣结果事实上都需要这些性质，而且在自然界中遇到的空间极少会因为不满足这两个假设之一而不是流形. 关于几个简单的例子，见习题 1-1 和 1-2；关于一个更复杂的例子（一个连通的、局部欧氏的、Hausdorff的但不是第二可数的空间），见 [LeeTM, 习题 4-6].
